\documentclass[aspectratio=169]{beamer}
\usetheme{Madrid}
\usecolortheme{seahorse}

\usepackage{graphicx}
\usepackage{booktabs}
\usepackage{amsmath}
\usepackage{amssymb}
\usepackage{tabularx}
\usepackage{array}
\usepackage{ragged2e}

\title[Chapter 14: Introduction to Inference]{\textbf{Chapter 14: Introduction to Inference}}
\subtitle{The Practice of Statistics in the Life Sciences \\ by Brigitte Baldi and David S. Moore}
\author{Instructor: Ali Raisolsadat}
\date{Part III: Foundations of Inference}

\setbeamertemplate{navigation symbols}{}
\setbeamertemplate{itemize items}[circle]
\setbeamertemplate{enumerate items}[default]
\graphicspath{{./}{figures/}}

\newcommand{\lectureimage}[2][0.90\linewidth]{%
  \IfFileExists{#2}{%
    \includegraphics[width=#1,height=0.74\textheight,keepaspectratio]{#2}%
  }{%
    \fbox{\parbox[c][0.38\textheight][c]{0.88\linewidth}{%
      \centering\small
      Figure file: \texttt{\detokenize{#2}}\\[0.15cm]
      Generate this figure by knitting the revised Chapter 14 RMarkdown notes.
    }}%
  }%
}

\begin{document}

% SLIDE 1: Title
\begin{frame}
  \titlepage
\end{frame}

% SLIDE 2: From Sampling Distributions to Inference
\begin{frame}{From Sampling Distributions to Inference}
Chapter 13 explained how sample statistics vary from one random sample to another.

\vspace{0.25cm}
Chapter 14 uses those sampling distributions to answer a new question:

\begin{center}
\large
\textbf{What can a sample tell us about an unknown population parameter?}
\end{center}

\vspace{0.2cm}
Statistical inference connects sample evidence to population conclusions while accounting for random sampling variation.
\end{frame}

% SLIDE 3: Two Main Goals of Inference
\begin{frame}{Two Main Goals of Inference}
Two major forms of inference are introduced in this chapter.

\begin{itemize}
  \item A \textbf{confidence interval} estimates an unknown population parameter and describes the precision of the estimate.
  \item A \textbf{hypothesis test} evaluates how compatible the observed data are with a specific claim about a population parameter.
\end{itemize}

\vspace{0.25cm}
Both methods depend on the sampling distributions developed in Chapter 13.
\end{frame}

% SLIDE 4: Simple Conditions for Introductory Inference
\begin{frame}{Simple Conditions for Introductory Inference}
The first procedures in this chapter use a simplified setting so that the reasoning is easy to see.

\begin{enumerate}
  \item The observations come from an appropriate random sampling process and can be treated as independent.
  \item The sampling distribution of $\bar{x}$ can be modelled by a Normal distribution.
  \item The population standard deviation $\sigma$ is treated as known.
\end{enumerate}

\vspace{0.2cm}
Later chapters replace the known value of $\sigma$ with the sample standard deviation $s$ and use Student's $t$ distribution.
\end{frame}

% SLIDE 5: Why Start with Known Sigma
\begin{frame}{Why Start with a Known Population Standard Deviation?}
Knowing the true value of $\sigma$ is unusual in real applications.

\vspace{0.2cm}
The assumption is useful here because it separates two ideas:

\begin{itemize}
  \item The \textbf{logic of inference} can be learned first.
  \item The additional mathematics required when $\sigma$ must be estimated can be introduced later.
\end{itemize}

\vspace{0.2cm}
The study design still matters. A correct formula cannot repair serious bias caused by an inappropriate sample.
\end{frame}

\section{Statistical Estimation}

% SLIDE 6: Statistical Estimation
\begin{frame}{Statistical Estimation}
A sample statistic can be used as an estimate of a population parameter.

\vspace{0.2cm}
For a population mean $\mu$, the natural point estimate is

\[
\bar{x}
\]

A point estimate gives one number. A confidence interval adds a margin of error so that sampling uncertainty is also described.
\end{frame}

% SLIDE 7: Height Example Question
\begin{frame}{Examples 14.1 and 14.2: Mean Height}
\begin{block}{Question}
A health survey provides a sample of $n=217$ eight-year-old boys from the population of interest. The sample mean height is $\bar{x}=132.5$ cm. For this introductory example, assume the population standard deviation is $\sigma=10$ cm.

How can the sample mean be used to estimate the unknown population mean height $\mu$ while accounting for random sampling variation?
\end{block}
\end{frame}

% SLIDE 8: Height Example Explanation Step 1
\begin{frame}{Examples 14.1 and 14.2: Sampling Distribution}
\begin{block}{Explanation}
The sampling distribution of $\bar{x}$ has centre $\mu$ and standard deviation
\[
\sigma_{\bar{x}}
=
\frac{\sigma}{\sqrt{n}}
\]
\end{block}

\pause

For this sample size,
\[
\sigma_{\bar{x}}
=
\frac{10}{\sqrt{217}}
\approx
0.679\text{ cm}
\]

\pause

This quantity is the \textbf{standard error} of the sample mean under the model.
\end{frame}

% SLIDE 9: Height Example Explanation Step 2
\begin{frame}{Examples 14.1 and 14.2: Approximate 95\% Distance}
\begin{block}{Explanation}
The 68--95--99.7 rule says that approximately 95\% of a Normal distribution lies within two standard deviations of its mean.
\end{block}

\pause

Using the exact standard error gives
\[
2(0.679)
\approx
1.36\text{ cm}
\]

\pause

The textbook-style calculation rounds the standard error to 0.7 cm first, which gives an approximate margin of error of
\[
2(0.7)
=
1.4\text{ cm}
\]
\end{frame}

% SLIDE 10: Height Example Solution
\begin{frame}{Examples 14.1 and 14.2: Confidence Interval}
\begin{block}{Solution}
The lower endpoint is
\[
132.5-1.4
=
131.1
\]
\end{block}

\pause

The upper endpoint is
\[
132.5+1.4
=
133.9
\]

\pause

The approximate 95\% confidence interval is
\[
(131.1,\ 133.9)\text{ cm}
\]

\pause

We are approximately 95\% confident that the population mean height lies between 131.1 cm and 133.9 cm because this method captures the true mean in about 95\% of repeated samples under the stated model.
\end{frame}

% SLIDE 11: Repeated Confidence Intervals Plot
\begin{frame}{Repeated 95\% Confidence Intervals}
\centering
\lectureimage[0.92\linewidth]{fig_ch14_sampling_dist.png}
\end{frame}

% SLIDE 12: Long-Run Coverage
\begin{frame}{What Does 95\% Coverage Mean?}
The population mean is fixed. The intervals change from sample to sample.


A 95\% confidence procedure has a \textbf{long-run coverage rate of 95\%} under the model.


This means that approximately 95\% of intervals produced by the method would contain the true parameter over many repeated random samples.


A small simulation does not need to show exactly 95\% coverage. The theoretical statement concerns the long-run procedure.
\end{frame}

\section{Margin of Error and Confidence Level}

% SLIDE 13: Anatomy of a Confidence Interval
\begin{frame}{Anatomy of a Confidence Interval}
Most confidence intervals have the form

\[
\text{estimate}
\pm
\text{margin of error}
\]

\vspace{0.3cm}
Three ideas determine the interval:

\begin{itemize}
  \item The \textbf{estimate} determines the centre.
  \item The \textbf{standard error} describes sampling variability.
  \item The \textbf{confidence level} determines how wide the interval must be for the desired long-run coverage.
\end{itemize}
\end{frame}

% SLIDE 14: The Estimate
\begin{frame}{The Estimate}
When estimating a population mean, the centre of the interval is usually the sample mean.

\[
\bar{x}
\]

For the height example,
\[
\bar{x}
=
132.5\text{ cm}
\]

The sample mean is unbiased for $\mu$ under the usual random-sampling model. This property does not guarantee that one particular sample mean is close to the population mean.
\end{frame}

% SLIDE 15: The Margin of Error
\begin{frame}{The Margin of Error}
For a confidence interval based on the Normal distribution,

\[
m
=
z^*
\frac{\sigma}{\sqrt{n}}
\]


The margin of error describes uncertainty caused by random sampling under the model.


It does not automatically account for other sources of error such as poor measurement, nonresponse or selection bias.
\end{frame}

% SLIDE 16: The Confidence Level
\begin{frame}{The Confidence Level}
The confidence level $C$ describes the long-run success rate of the interval-building procedure.

\begin{itemize}
  \item A 90\% interval uses a smaller critical value than a 95\% interval.
  \item A 99\% interval uses a larger critical value than a 95\% interval.
\end{itemize}


Higher confidence requires a wider interval when the sample size and population standard deviation remain unchanged.
\end{frame}

% SLIDE 17: Precision and Confidence
\begin{frame}{Precision and Confidence}
Two changes have predictable effects on interval width.


\begin{block}{Larger Sample Size}
A larger sample reduces the standard error because
\[
\frac{\sigma}{\sqrt{n}}
\]
becomes smaller as $n$ increases.
\end{block}


\begin{block}{Higher Confidence Level}
A higher confidence level uses a larger $z^*$ value, which makes the interval wider.
\end{block}
\end{frame}

\section{Interpreting Confidence Intervals}

% SLIDE 18: Correct Interpretation
\begin{frame}{How to Interpret a Confidence Interval}
A suitable interpretation is:

\begin{block}{95\% Confidence}
We are 95\% confident that the population parameter lies between the reported endpoints because the procedure used to construct the interval captures the true parameter in 95\% of repeated samples under the stated conditions.
\end{block}

The confidence level describes the reliability of the \textbf{method} over repeated sampling.
\end{frame}

% SLIDE 19: What 95 Percent Confidence Does Not Mean
\begin{frame}{What 95\% Confidence Does Not Mean}
After an interval has been calculated, the population parameter is fixed.


The interval either contains the parameter or it does not. The difficulty is that the parameter is unknown.


For a frequentist confidence interval, we do not say that there is a 95\% probability that the fixed parameter lies inside the specific interval that was already calculated.
\end{frame}

% SLIDE 20: Twenty-Five Confidence Intervals Plot
\begin{frame}{Twenty-Five Simulated 95\% Confidence Intervals}
\centering
\lectureimage[0.92\linewidth]{fig_ch14_25_intervals.png}
\end{frame}

\section{Confidence Intervals for the Mean}

% SLIDE 21: Common Critical Values
\begin{frame}{Common Critical Values}
The critical value $z^*$ is chosen so that the central area under the Standard Normal curve equals the desired confidence level.

\vspace{0.2cm}
\begin{center}
\begin{tabular}{cc}
\toprule
\textbf{Confidence Level} & \textbf{Critical Value $z^*$} \\
\midrule
90\% & 1.645 \\
95\% & 1.960 \\
99\% & 2.576 \\
\bottomrule
\end{tabular}
\end{center}
\end{frame}

% SLIDE 22: Confidence Interval Formula
\begin{frame}{Confidence Interval for a Population Mean}
Under the introductory conditions with known $\sigma$,

\[
\bar{x}
\pm
z^*
\frac{\sigma}{\sqrt{n}}
\]


The same formula can be viewed as

\[
\text{estimate}
\pm
\text{critical value}\times\text{standard error}
\]


The estimate determines the centre. The standard error measures sampling variation. The critical value sets the desired confidence level.
\end{frame}

% SLIDE 23: Four-Step Process for Confidence Intervals
\begin{frame}{The Four-Step Process for Confidence Intervals}
A consistent process keeps the calculation connected to the scientific question.

\begin{enumerate}
  \item \textbf{STATE:} The population question and parameter are identified.
  \item \textbf{PLAN:} The confidence level and appropriate procedure are selected.
  \item \textbf{SOLVE:} The conditions are assessed before the interval is calculated.
  \item \textbf{CONCLUDE:} The interval is interpreted in the context of the population.
\end{enumerate}
\end{frame}

% SLIDE 24: Leaf Surface Area Question
\begin{frame}{Example 14.3: Leaf Surface Area}
\begin{block}{Question}
A botanist wants to estimate the mean surface area of mature \textit{Monstera deliciosa} leaves in a greenhouse population. A random sample contains 31 mature leaves. The population standard deviation is treated as known at $\sigma=15$ cm$^2$.

What is a 95\% confidence interval for the population mean leaf surface area?
\end{block}
\end{frame}

% SLIDE 25: Leaf Surface Area Data
\begin{frame}{Example 14.3: Observed Leaf Data}
\scriptsize
\begin{center}
\begin{tabular}{rrrrrrrr}
114 & 100 & 104 & 89 & 102 & 91 & 114 & 114 \\
103 & 105 & 108 & 130 & 120 & 132 & 111 & 128 \\
118 & 119 & 86 & 72 & 111 & 103 & 74 & 112 \\
107 & 103 & 98 & 96 & 112 & 112 & 93 & \\
\end{tabular}
\end{center}
\normalsize

\pause

\begin{block}{Explanation}
The parameter is the population mean leaf surface area $\mu$. The first step is to assess whether the sampling model is reasonable before calculating the interval.
\end{block}
\end{frame}

% SLIDE 26: Leaf Surface Area Plot
\begin{frame}{Example 14.3: Leaf Surface Area Distribution}
\centering
\lectureimage[0.92\linewidth]{fig_ch14_leaf_dotplot.png}
\end{frame}

% SLIDE 27: Leaf Surface Area Conditions
\begin{frame}{Example 14.3: Checking the Conditions}
\begin{block}{Explanation}
The sample was described as randomly selected, which supports the sampling condition.
\end{block}

\pause

The observed distribution is somewhat left-skewed with no extreme outliers.

\pause

With $n=31$, the Central Limit Theorem provides additional support for an approximately Normal sampling distribution of $\bar{x}$ when the population has finite variance.

\pause

The population standard deviation is treated as known at
\[
\sigma=15\text{ cm}^2
\]
\end{frame}

% SLIDE 28: Leaf Surface Area Calculation
\begin{frame}{Example 14.3: Calculating the Interval}
The sample mean is
\[
\bar{x}
\approx
105.84
\]

\pause

The standard error is
\[
SE
=
\frac{15}{\sqrt{31}}
\approx
2.69
\]

\pause

The 95\% critical value is
\[
z^*=1.960
\]

\pause

The margin of error is
\[
m
=
1.960(2.69)
\approx
5.28
\]
\end{frame}

% SLIDE 29: Leaf Surface Area Solution
\begin{frame}{Example 14.3: Confidence Interval Solution}
\begin{block}{Solution}
The interval is
\[
105.84
\pm
5.28
\]
\end{block}

\pause

Therefore,
\[
(100.56,\ 111.12)
\]

\pause

We are 95\% confident that the population mean surface area of mature leaves in the greenhouse population is between approximately 100.6 cm$^2$ and 111.1 cm$^2$.
\end{frame}

% SLIDE 30: Leaf Surface Area R Code
\begin{frame}[fragile]{Example 14.3: Confidence Interval in R}
\scriptsize
\begin{verbatim}
# The 31 observed leaf surface areas are stored in a vector.
leaf_data <- c(
  114, 100, 104, 89, 102, 91, 114, 114,
  103, 105, 108, 130, 120, 132, 111, 128,
  118, 119, 86, 72, 111, 103, 74, 112,
  107, 103, 98, 96, 112, 112, 93
)

# The population standard deviation is treated as known.
sigma_leaf <- 15
xbar_leaf <- mean(leaf_data)
n_leaf <- length(leaf_data)

# qnorm(0.975) gives the 95% Normal critical value.
z_star_leaf <- qnorm(0.975)
se_leaf <- sigma_leaf / sqrt(n_leaf)
margin_leaf <- z_star_leaf * se_leaf

c(xbar_leaf - margin_leaf, xbar_leaf + margin_leaf)
\end{verbatim}
\end{frame}

% SLIDE 31: Beyond the Simple Conditions
\begin{frame}{Looking Beyond the Simple Conditions}
The known-$\sigma$ assumption is mainly a teaching device.

\vspace{0.2cm}
In most real studies, the population standard deviation is unknown.


Later chapters replace $\sigma$ with the sample standard deviation $s$ and use Student's $t$ distribution.


The central logic remains the same: an estimate is combined with a standard error and an appropriate critical value.
\end{frame}

\section{Hypothesis Testing}

% SLIDE 32: From Estimation to Testing
\begin{frame}{From Estimation to Hypothesis Testing}
A confidence interval begins with an unknown parameter and asks for a plausible range of values.


A hypothesis test begins with a specific population claim and asks whether the observed sample data are compatible with that claim.


The central question is:

\begin{center}
\large
\textbf{If the null hypothesis were true, how unusual would the observed sample result be?}
\end{center}
\end{frame}

% SLIDE 33: Phosphorus Example Question
\begin{frame}{Examples 14.4 and 14.5: Phosphorus Levels}
\begin{block}{Question}
In the general adult population, blood phosphorus levels are modelled as Normal with mean 1.2 mmol/L and standard deviation 0.1 mmol/L. Researchers suspect that the mean may be lower among adults aged 75 to 79.

A retrospective chart review provides measurements from 12 patients. Does the sample provide evidence that the population mean for this older group is below 1.2 mmol/L?
\end{block}
\end{frame}

% SLIDE 34: Phosphorus Data
\begin{frame}{Examples 14.4 and 14.5: Phosphorus Data}
\begin{center}
\begin{tabular}{rrrrrr}
\toprule
1.26 & 1.39 & 1.00 & 1.00 & 1.00 & 1.10 \\
0.87 & 1.23 & 1.19 & 1.29 & 1.03 & 1.18 \\
\bottomrule
\end{tabular}
\end{center}

\pause

The sample mean is
\[
\bar{x}
\approx
1.1283\text{ mmol/L}
\]

\pause

The sample mean is below 1.2, but sampling variation alone can produce differences between $\bar{x}$ and $\mu$.
\end{frame}

% SLIDE 35: Phosphorus Hypotheses
\begin{frame}{Examples 14.4 and 14.5: Hypotheses}
\begin{block}{Explanation}
The test compares the observed difference with the variation expected if the general-adult mean still applied to the older population.
\end{block}

\pause

\begin{block}{Solution}
The null hypothesis is
\[
H_0:\mu=1.2
\]
\end{block}

\pause

The alternative hypothesis is
\[
H_a:\mu<1.2
\]

The alternative is one-sided because the scientific question concerns a decrease.
\end{frame}

% SLIDE 36: Null and Alternative Hypotheses
\begin{frame}{The Null and Alternative Hypotheses}
The \textbf{null hypothesis} gives the specific population value or model being evaluated.

\[
H_0:\mu=\mu_0
\]


The \textbf{alternative hypothesis} gives the competing population claim that is supported when the data provide sufficient evidence against $H_0$.


Hypotheses refer to population parameters. They do not describe sample statistics that have already been observed.
\end{frame}

% SLIDE 37: One-Sided and Two-Sided Alternatives
\begin{frame}{One-Sided and Two-Sided Alternatives}
A lower-tail alternative has the form
\[
H_a:\mu<\mu_0
\]

An upper-tail alternative has the form
\[
H_a:\mu>\mu_0
\]

A two-sided alternative has the form
\[
H_a:\mu\ne\mu_0
\]

The direction should be determined by the scientific question before the sample result is used to choose the test.
\end{frame}

% SLIDE 38: Aspirin Example Question
\begin{frame}{Example 14.6: Aspirin Tablet Manufacturing}
\begin{block}{Question}
A pharmaceutical process is designed to produce aspirin tablets with a population mean active-ingredient amount of 325 mg. Both underfilling and overfilling are undesirable.

What hypotheses should be used to evaluate whether the population mean differs from the target?
\end{block}
\end{frame}

% SLIDE 39: Aspirin Example Solution
\begin{frame}{Example 14.6: Choosing the Hypotheses}
\begin{block}{Explanation}
The reference value is 325 mg. The scientific concern includes departures above or below the target.
\end{block}

\pause

\begin{block}{Solution}
The null hypothesis is
\[
H_0:\mu=325
\]
\end{block}

\pause

The alternative hypothesis is
\[
H_a:\mu\ne325
\]

This is a two-sided test because deviations in either direction matter.
\end{frame}

\section{P-Values and Statistical Significance}

% SLIDE 40: Logic of the P-Value
\begin{frame}{The Logic of a P-Value}
A hypothesis test begins by temporarily assuming that $H_0$ is true.


The null sampling distribution then tells us which sample results are common and which are unusual.


A result that is very unusual under the null model provides evidence against that model.
\end{frame}

% SLIDE 41: The Test Statistic
\begin{frame}{The Test Statistic}
For a one-sample $z$ test of a population mean,

\[
z
=
\frac{\bar{x}-\mu_0}{\sigma/\sqrt{n}}
\]


The numerator measures the observed difference from the null value.


The denominator is the standard error expected from random sampling.


The resulting $z$ value tells us how many standard errors the observed sample mean is from the null mean.
\end{frame}

% SLIDE 42: The P-Value
\begin{frame}{The P-Value}
\begin{block}{Definition}
The P-value is the probability, calculated under the assumption that the null hypothesis and model conditions are true, of obtaining a test statistic at least as extreme as the observed statistic in the direction or directions specified by the alternative hypothesis.
\end{block}


A P-value is not the probability that $H_0$ is true.


A P-value is not the probability that the observed result happened only because of chance.
\end{frame}

% SLIDE 43: Interpreting P-Values
\begin{frame}{Interpreting P-Values}
\begin{itemize}
  \item A small P-value provides stronger evidence against $H_0$ because the observed result is unusual under the null model.
  \item A large P-value means the observed result remains reasonably compatible with $H_0$.
  \item Failure to reject $H_0$ does not prove that $H_0$ is true.
\end{itemize}

\vspace{0.2cm}
Statistical significance does not measure the practical or biological importance of an effect.
\end{frame}

% SLIDE 44: Phosphorus P-Value Question
\begin{frame}{Example 14.7: One-Sided Phosphorus Test}
\begin{block}{Question}
For the phosphorus study,
\[
H_0:\mu=1.2
\]
\[
H_a:\mu<1.2
\]
with $n=12$, $\bar{x}\approx1.1283$ mmol/L and known $\sigma=0.1$ mmol/L.

What is the one-sided P-value?
\end{block}
\end{frame}

% SLIDE 45: Phosphorus Standard Error
\begin{frame}{Example 14.7: Standard Error}
\begin{block}{Explanation}
Under $H_0$, the sample mean has centre 1.2 mmol/L.
\end{block}

\pause

The standard error is
\[
SE
=
\frac{0.1}{\sqrt{12}}
\approx
0.0289
\]

\pause

The observed sample mean is lower than the null mean, so the next step is to measure that difference in standard-error units.
\end{frame}

% SLIDE 46: Phosphorus Test Statistic
\begin{frame}{Example 14.7: Test Statistic and Tail}
The test statistic is
\[
z
=
\frac{1.1283-1.2}{0.0289}
\approx
-2.48
\]

\pause

The alternative hypothesis is
\[
H_a:\mu<1.2
\]

\pause

Therefore, the P-value is the Standard Normal area to the \textbf{left} of the observed $z$ statistic.
\end{frame}

% SLIDE 47: Phosphorus P-Value Solution
\begin{frame}{Example 14.7: P-Value Solution}
\begin{block}{Solution}
\[
P(Z\le-2.48)
\approx
0.0065
\]
\end{block}

\pause

Under the null model, a sample mean of approximately 1.1283 mmol/L or lower would occur in about 0.65\% of repeated samples of size 12.

\pause

The data provide strong evidence against $H_0$ under the stated assumptions and support the conclusion that the population mean is below 1.2 mmol/L.
\end{frame}

% SLIDE 48: Phosphorus R Code
\begin{frame}[fragile]{Example 14.7: One-Sided Test in R}
\scriptsize
\begin{verbatim}
# The observed phosphorus measurements are stored in mmol/L.
phosphorus <- c(
  1.26, 1.39, 1.00, 1.00, 1.00, 1.10,
  0.87, 1.23, 1.19, 1.29, 1.03, 1.18
)

# The null mean and population standard deviation define the model.
mu_null_phos <- 1.2
sigma_phos <- 0.1

xbar_phos <- mean(phosphorus)
n_phos <- length(phosphorus)
se_phos <- sigma_phos / sqrt(n_phos)
z_phos <- (xbar_phos - mu_null_phos) / se_phos

# The alternative is lower-tailed, so pnorm() gives the left-tail area.
p_phos <- pnorm(z_phos)
p_phos
\end{verbatim}
\end{frame}

% SLIDE 49: Aspirin P-Value Question
\begin{frame}{Example 14.8: Two-Sided Aspirin Test}
\begin{block}{Question}
A manufacturing process targets $\mu=325$ mg with known $\sigma=5$ mg. A random sample of $n=10$ tablets has mean $\bar{x}=326.9$ mg.

The hypotheses are
\[
H_0:\mu=325
\]
\[
H_a:\mu\ne325
\]

What is the two-sided P-value?
\end{block}
\end{frame}

% SLIDE 50: Aspirin Explanation
\begin{frame}{Example 14.8: Standard Error and Test Statistic}
\begin{block}{Explanation}
The standard error is
\[
SE
=
\frac{5}{\sqrt{10}}
\approx
1.581
\]
\end{block}

\pause

The test statistic is
\[
z
=
\frac{326.9-325}{1.581}
\approx
1.202
\]

\pause

The two-sided alternative requires equally extreme outcomes in both tails.
\end{frame}

% SLIDE 51: Aspirin Solution
\begin{frame}{Example 14.8: Two-Sided P-Value Solution}
\begin{block}{Solution}
For the symmetric Standard Normal distribution,
\[
2P(Z\ge1.202)
\approx
0.2295
\]
\end{block}

\pause

A P-value of about 0.2295 is not small.

\pause

The data do not provide statistically significant evidence that the population mean differs from 325 mg. This conclusion does not prove that the process mean is exactly 325 mg.
\end{frame}

% SLIDE 52: Aspirin R Code
\begin{frame}[fragile]{Example 14.8: Two-Sided Test in R}
\scriptsize
\begin{verbatim}
# The null value is the target tablet amount.
mu_null_aspirin <- 325
sigma_aspirin <- 5
n_aspirin <- 10
xbar_aspirin <- 326.9

# The z statistic measures the standardized distance from the target.
se_aspirin <- sigma_aspirin / sqrt(n_aspirin)
z_aspirin <- (xbar_aspirin - mu_null_aspirin) / se_aspirin

# abs() lets the same two-sided formula work for positive or negative z.
p_aspirin <- 2 * pnorm(-abs(z_aspirin))
p_aspirin
\end{verbatim}
\end{frame}

\section{Significance Levels and One-Sample Tests}

% SLIDE 53: Significance Levels
\begin{frame}{Significance Levels $\alpha$}
The significance level is a decision threshold chosen before the analysis.


The usual decision rule is

\begin{itemize}
  \item When $P\le\alpha$, the result is statistically significant and $H_0$ is rejected.
  \item When $P>\alpha$, the result is not statistically significant and $H_0$ is not rejected.
\end{itemize}


Common values include
\[
\alpha=0.05
\]
\[
\alpha=0.01
\]
\end{frame}

% SLIDE 54: What Alpha Means
\begin{frame}{What Does $\alpha$ Mean?}
For a correctly calibrated test, $\alpha$ controls the long-run probability of a Type I error when the null hypothesis is true.


The statement
\[
\alpha=0.05
\]
does not mean that there is a 5\% probability that $H_0$ is true.


The word \textbf{significant} also does not mean that an effect is large or biologically important.
\end{frame}

% SLIDE 55: One-Sample Z Test
\begin{frame}{The One-Sample $z$ Test for a Mean}
Under the introductory conditions, the null hypothesis is
\[
H_0:\mu=\mu_0
\]

The test statistic is
\[
z
=
\frac{\bar{x}-\mu_0}{\sigma/\sqrt{n}}
\]

The test statistic measures the observed difference from the null value in units of the standard error.
\end{frame}

% SLIDE 56: Which Tail Gives the P-Value
\begin{frame}{Which Tail Gives the P-Value?}
The alternative hypothesis determines the relevant tail or tails.

\begin{itemize}
  \item For $H_a:\mu>\mu_0$, the P-value is the area to the right of the observed $z$ value.
  \item For $H_a:\mu<\mu_0$, the P-value is the area to the left of the observed $z$ value.
  \item For $H_a:\mu\ne\mu_0$, the P-value includes equally extreme outcomes in both tails.
\end{itemize}


For a two-sided $z$ test in R,
\[
P=2\,\texttt{pnorm}(-|z|)
\]
\end{frame}

\section{Body Temperature Example}

% SLIDE 57: Body Temperature Question
\begin{frame}{Example 14.9: Body Temperature}
\begin{block}{Question}
An historical reference value for mean human body temperature is 98.6$^\circ$F. A later study reports oral temperatures from $n=130$ healthy adult volunteers with sample mean $\bar{x}=98.25^\circ$F.

For this introductory exercise, assume the observations can be treated as independent and representative of the population being studied. Also assume known $\sigma=0.6^\circ$F.

Do the data provide evidence that the population mean differs from 98.6$^\circ$F?
\end{block}
\end{frame}

% SLIDE 58: Body Temperature Hypotheses and Standard Error
\begin{frame}{Example 14.9: Plan and Standard Error}
\begin{block}{Explanation}
The question concerns a difference in either direction, so a two-sided test is appropriate.
\end{block}

\pause

\[
H_0:\mu=98.6
\]
\[
H_a:\mu\ne98.6
\]

\pause

The standard error is
\[
SE
=
\frac{0.6}{\sqrt{130}}
\approx
0.0526
\]
\end{frame}

% SLIDE 59: Body Temperature Test Statistic
\begin{frame}{Example 14.9: Test Statistic}
The test statistic is
\[
z
=
\frac{98.25-98.6}{0.0526}
\approx
-6.65
\]

\pause

The observed sample mean is about 6.65 standard errors below the null mean.

\pause

A result at least this far from 98.6 in either direction would be extremely unusual under the null model.
\end{frame}

% SLIDE 60: Body Temperature Solution
\begin{frame}{Example 14.9: P-Value and Conclusion}
\begin{block}{Solution}
The two-sided P-value is approximately
\[
2.9\times10^{-11}
\]
\end{block}

\pause

The P-value is extremely small. Under the stated assumptions, the data provide extremely strong evidence against $H_0$.

\pause

The result is not impossible under $H_0$. Its probability is simply extremely small. The strength of the population conclusion also depends on whether the study sample appropriately represents the population of interest.
\end{frame}

% SLIDE 61: Body Temperature R Code
\begin{frame}[fragile]{Example 14.9: Body Temperature Test in R}
\scriptsize
\begin{verbatim}
# The null model uses the historical reference mean.
mu_null_temp <- 98.6
xbar_temp <- 98.25
n_temp <- 130
sigma_temp <- 0.6

# The standard error and z statistic describe the observed difference.
se_temp <- sigma_temp / sqrt(n_temp)
z_temp <- (xbar_temp - mu_null_temp) / se_temp

# The alternative is two-sided, so both tails are included.
p_temp <- 2 * pnorm(-abs(z_temp))

c(Standard_Error = se_temp, Z_Statistic = z_temp, P_Value = p_temp)
\end{verbatim}
\end{frame}

\section{Tests from Confidence Intervals}

% SLIDE 62: Confidence Interval and Test Connection
\begin{frame}{The Confidence Interval and Test Connection}
Two-sided confidence intervals and two-sided hypothesis tests are closely connected when they use the same model.


A test at significance level $\alpha$ corresponds to a confidence level
\[
C=1-\alpha
\]


For example,
\[
\alpha=0.05
\]
corresponds to
\[
C=0.95
\]
\end{frame}

% SLIDE 63: Decision from a Confidence Interval
\begin{frame}{Using a Confidence Interval to Test a Null Value}
Under the same two-sided model,

\begin{itemize}
  \item A null value outside the corresponding confidence interval leads to rejection of $H_0$.
  \item A null value inside the corresponding confidence interval leads to failure to reject $H_0$.
\end{itemize}

\vspace{0.2cm}
This equivalence requires the same standard error, significance level and two-sided model.
\end{frame}

% SLIDE 64: Body Temperature Confidence Interval Question
\begin{frame}{Example 14.10: Body Temperature Confidence Interval}
\begin{block}{Question}
Using $\bar{x}=98.25^\circ$F, $n=130$ and known $\sigma=0.6^\circ$F, what is the 95\% confidence interval for the population mean?

How does the interval lead to the same conclusion as the two-sided test of
\[
H_0:\mu=98.6
\]
at $\alpha=0.05$?
\end{block}
\end{frame}

% SLIDE 65: Body Temperature Confidence Interval Explanation
\begin{frame}{Example 14.10: Calculating the Interval}
\begin{block}{Explanation}
The 95\% critical value is
\[
z^*=1.960
\]
\end{block}

\pause

The standard error is
\[
SE
\approx
0.0526
\]

\pause

The margin of error is
\[
m
=
1.960(0.0526)
\approx
0.1031
\]
\end{frame}

% SLIDE 66: Body Temperature Confidence Interval Solution
\begin{frame}{Example 14.10: Confidence Interval Solution}
\begin{block}{Solution}
The interval is
\[
98.25
\pm
0.1031
\]
\end{block}

\pause

Therefore,
\[
(98.15,\ 98.35)
\]

\pause

The null value 98.6$^\circ$F lies outside the 95\% confidence interval. The corresponding two-sided test rejects $H_0$ at the 5\% significance level.
\end{frame}

% SLIDE 67: Body Temperature Confidence Interval R Code
\begin{frame}[fragile]{Example 14.10: Confidence Interval in R}
\scriptsize
\begin{verbatim}
# qnorm(0.975) gives the critical value for a 95% interval.
z_star_temp <- qnorm(0.975)

# The confidence interval uses the same standard error as the test.
margin_temp <- z_star_temp * se_temp
lower_temp <- xbar_temp - margin_temp
upper_temp <- xbar_temp + margin_temp

round(
  c(Lower_Bound = lower_temp, Upper_Bound = upper_temp),
  2
)
\end{verbatim}
\end{frame}

\section{Chapter Summary}

% SLIDE 68: Chapter Summary
\begin{frame}{Chapter 14 Summary}
\begin{itemize}
  \item Statistical inference uses sample data to learn about population parameters while accounting for sampling variation.
  \item Confidence intervals combine an estimate with a margin of error.
  \item A P-value measures how unusual the observed evidence would be under the null model.
  \item The significance level $\alpha$ provides a decision threshold chosen before the analysis.
  \item Two-sided confidence intervals and two-sided tests are directly connected when they use the same statistical model.
\end{itemize}
\end{frame}

% SLIDE 69: Consistent Inference Workflow
\begin{frame}{A Consistent Inference Workflow}
A first-year inference problem becomes easier when the same reasoning is used each time.

\begin{enumerate}
  \item The population parameter and scientific question are identified.
  \item The study design and assumptions are assessed.
  \item The sampling distribution and standard error are identified.
  \item The confidence interval or test is calculated.
  \item The numerical result is interpreted in context.
  \item The conclusion respects both statistical evidence and the limitations of the study design.
\end{enumerate}
\end{frame}

\end{document}
